% ===========================================================================
%  Chapter 03 homework — Rings
%  Garling, A Course in Galois Theory
%
%  DIVISION OF LABOUR (see AI_INSTRUCTIONS.md §7):
%    The assistant transcribes each assigned problem statement faithfully and
%    emits an empty, marked solution region beneath it.
%    The user types the mathematics inside those regions. No assistant writes
%    inside a SOLUTION region in normal mode.
% ===========================================================================
\documentclass[11pt]{article}
\usepackage{coursemacros}

\title{Chapter 03 Homework --- Rings}
\author{Mikey Ferguson}
\date{\today}

\begin{document}
\maketitle

% ---------------------------------------------------------------------------
% Assigned set (2026-09-02), in the order given by the student's own sheet:
%   3.1, 3.2, 3.3, 3.8, 3.10, 3.11, 3.12, 3.14, 3.15, 3.17, 3.25, 3.26
% Statements are transcribed as each problem is reached; solution regions
% stay empty until the student's answer is agreed correct.
% ---------------------------------------------------------------------------

\begin{problem}{03.1}
  Suppose that $S$ is a set and $R$ is a ring. Let $R^{S}$ denote the set of
  all mappings from $S$ to $R$. Show that $R^{S}$ is a ring, under the
  operations defined by
  \[
    (f+g)(s) = f(s) + g(s), \qquad (fg)(s) = f(s)g(s).
  \]
  Show that if $S$ has more than one element then there exist non-zero
  elements $f$ and $g$ in $R^{S}$ for which $fg = 0$.
\end{problem}
% ===== SOLUTION 03.1 =====
% Transcribed from the student's slate, 2026-09-02 (attempts 1-6, turns
% t0002-t0004; pages filed in ../handwritten/20260902-3.1-attempt*.png).
% Typeset exactly as written, on their instruction ("What I wrote is good
% enough"). Faithful transcription: same steps, same order, nothing added.
% Not covered by what they wrote, and flagged to them on card 0009:
%   - only distributivity is verified; the remaining ring axioms are asserted;
%   - no line asserting f and g are non-zero (needs 1_R =/= 0_R);
%   - no case split showing (fg)(s) = 0_R at every s.
\begin{solution}
  Write $f_{1}, f_{0} \in R^{S}$ for the functions
  \[
    f_{1}(s) = 1_{R}, \qquad f_{0}(s) = 0_{R},
  \]
  and clearly, by the specified function properties, we have all ring
  properties.

  Prove $(f+g)h = fh + gh$. Let $s \in S$. Show
  $\bigl((f+g)h\bigr)(s) = (fh)(s) + (gh)(s)$:
  \[
    \bigl((f+g)h\bigr)(s)
      = \bigl(f(s) + g(s)\bigr)h(s)
      = f(s)h(s) + g(s)h(s)
      = (fh)(s) + (gh)(s).
  \]

  If $\abs{S} > 1$, there exist $s_{1}, s_{2} \in S$ where $s_{1} \neq s_{2}$.
  Let
  \[
    f(s) = \begin{cases} 1_{R} & s = s_{1} \\ 0_{R} & \text{else} \end{cases}
    \qquad
    g(s) = \begin{cases} 1_{R} & s = s_{2} \\ 0_{R} & \text{else} \end{cases}
  \]
  Then $f(s)g(s) = 0$, so they are zero divisors.
\end{solution}
% ===== END SOLUTION 03.1 =====

\begin{problem}{03.2}
  Suppose that $R$ is an integral domain, with field of fractions $F$. Show
  that the field of fractions of $R[x_{1},\ldots,x_{n}]$ can be identified
  naturally with $F(x_{1},\ldots,x_{n})$.
\end{problem}
% ===== SOLUTION 03.2 =====
% TODO(mferguson): your work goes here.
% ===== END SOLUTION 03.2 =====

\begin{problem}{03.3}
  \todo{statement not yet transcribed}
\end{problem}
% ===== SOLUTION 03.3 =====
% TODO(mferguson): your work goes here.
% ===== END SOLUTION 03.3 =====

\begin{problem}{03.8}
  Suppose that $K$ is a field. If $f = a_{0} + a_{1}x + \cdots + a_{n}x^{n} \in
  K[x]$ and $k \in K$, let $(\Phi(f))(k) = a_{0} + a_{1}k + \cdots + a_{n}k^{n}$.
  Show that $\Phi$ is a ring homomorphism from $K[x]$ to $K^{K}$. Show that if
  $K$ is finite then $\Phi$ is an epimorphism, but not a monomorphism. What
  happens if $K$ is infinite?
\end{problem}
% ===== SOLUTION 03.8 =====
% PARTIAL. Transcribed from the student's slate, 2026-09-02, as each piece was
% agreed correct. Faithful transcription: same steps, same order, nothing added.
%   - additive half of the homomorphism clause: turn t0011 rev 5, page filed at
%     ../handwritten/20260902-3.8-attempt5-additive-correct.png (card 0037);
%   - multiplicative half: turn t0011 rev 12, page filed at
%     ../handwritten/20260902-3.8-attempt12-product-correct.png (card 0043);
%   - not a monomorphism, for K = F_2: turn t0011 rev 3, page filed at
%     ../handwritten/20260902-3.8-attempt3-kernel.png (card 0035);
%   - epimorphism for finite K: turn t0011 rev 19, page filed at
%     ../handwritten/20260903-3.8-attempt19-epimorphism-correct.png (card 0050);
%   - not a monomorphism, for general finite K: turn t0011 rev 20, page filed at
%     ../handwritten/20260903-3.8-attempt20-kernel-general-correct.png (card
%     0051). Their page exhibits f = prod_{k in K} delta_k and nothing else, so
%     that is all that is set. Not supplied here, and flagged on card 0051:
%       - no line showing f(c) = 0 for arbitrary c in K (some k =/= c kills it);
%       - no line showing f =/= 0, which needs K[x] to have no zero divisors.
% Two departures from the ink in the multiplicative half, both recorded rather
% than silently repaired:
%   - they wrote the product in K^K as "o" (composition); it is the pointwise
%     product of 3.1, so it is set here as juxtaposition. Notation, not a step.
%   - their page has no closing "as k was arbitrary" line, unlike the additive
%     half. Not added. Still owed.
% Left implicit in the epimorphism half and not supplied here:
%   - no general definition of delta_c is on any of their pages; they built it
%     by hand over F_3 and F_5 only. The general product form is still owed;
%   - no line saying the sum is finite because K is;
%   - the discarded terms are h(k)*0; the collapse is shown, not justified;
%   - no closing "as k_1 was arbitrary" line.
% Still owed by the student, regions left out of the write-up until then:
%   the "as k was arbitrary" line for the product; the general delta_c;
%   non-injectivity for general finite K; the infinite case.
\begin{solution}
  Let $f = \sum_{i} a_{i}x^{i}$ and $g = \sum_{i} b_{i}x^{i}$ in $K[x]$.

  \emph{Additive half.} Let $k \in K$. Then
  \[
    \bigl(\Phi(f+g)\bigr)(k)
      = \Bigl(\Phi\bigl(\textstyle\sum_{i}(a_{i}+b_{i})x^{i}\bigr)\Bigr)(k)
      = \sum_{i}(a_{i}+b_{i})k^{i}
      = \sum_{i} a_{i}k^{i} + \sum_{i} b_{i}k^{i}
  \]
  \[
      = \Bigl(\Phi\bigl(\textstyle\sum_{i} a_{i}x^{i}\bigr)\Bigr)(k)
        + \Bigl(\Phi\bigl(\textstyle\sum_{i} b_{i}x^{i}\bigr)\Bigr)(k)
      = \bigl(\Phi(f)\bigr)(k) + \bigl(\Phi(g)\bigr)(k).
  \]
  As $k \in K$ was arbitrary, $\Phi(f+g) = \Phi(f) + \Phi(g)$ in $K^{K}$.

  \emph{Multiplicative half.} Let $k \in K$. Then
  \[
    \bigl(\Phi(fg)\bigr)(k)
      = \Bigl(\Phi\Bigl(\textstyle\sum_{m}
          \bigl(\sum_{i+j=m} a_{i}b_{j}\bigr)x^{m}\Bigr)\Bigr)(k)
      = \sum_{m}\Bigl(\sum_{i+j=m} a_{i}b_{j}\Bigr)k^{m}
  \]
  \[
      = \Bigl(\sum_{i} a_{i}k^{i}\Bigr)\Bigl(\sum_{j} b_{j}k^{j}\Bigr)
      = \Bigl(\Phi\bigl(\textstyle\sum_{i} a_{i}x^{i}\bigr)\Bigr)(k)
        \cdot \Bigl(\Phi\bigl(\textstyle\sum_{j} b_{j}x^{j}\bigr)\Bigr)(k)
      = \bigl(\Phi(f)\bigr)(k)\bigl(\Phi(g)\bigr)(k).
  \]

  \emph{$\Phi$ is an epimorphism when $K$ is finite.} Let $K$ be a finite field
  and let $h \in K^{K}$ be arbitrary. We find $f \in K[x]$ with
  $\bigl(\Phi(f)\bigr)(k) = h(k)$. Take
  \[
    f(x) = \sum_{k} h(k)\,\delta_{k}(x).
  \]
  Let $k_{1} \in K$. Then
  \[
    \bigl(\Phi(f)\bigr)(k_{1})
      = \Bigl(\Phi\Bigl(\textstyle\sum_{k} h(k)\delta_{k}(x)\Bigr)\Bigr)(k_{1})
      = \sum_{k} h(k)\,\delta_{k}(k_{1})
      = h(k_{1})\,\delta_{k_{1}}(k_{1})
      = h(k_{1}) \cdot 1
      = h(k_{1}).
  \]
  So $\Phi$ is an epimorphism.

  \emph{$\Phi$ is not a monomorphism.} Take $K = \Fq{2}$, $f = x^{2} + x$ and
  $g = 0$. Then $f \neq g$ in $\Fq{2}[x]$, but
  \[
    \bigl(\Phi(f)\bigr)(0) = 0 = \bigl(\Phi(g)\bigr)(0),
    \qquad
    \bigl(\Phi(f)\bigr)(1) = 1 + 1 = 0 = \bigl(\Phi(g)\bigr)(1),
  \]
  so $\Phi(f) = \Phi(g)$ in $\Fq{2}^{\Fq{2}}$ and $\Ker \Phi \neq \{0\}$.

  For an arbitrary finite field $K$, take
  \[
    f(x) = \prod_{k \in K} \delta_{k}(x),
  \]
  a non-zero element of $\Ker \Phi$.
\end{solution}
% ===== END SOLUTION 03.8 =====

\begin{problem}{03.10}
  What are the units in $R[x]$, where $R$ is an integral domain? What are the
  units in $\mathbb{Z}_{4}[x]$?
\end{problem}
% ===== SOLUTION 03.10 =====
% PARTIAL. First question only (units in R[x] for R an integral domain),
% transcribed from the student's slate, 2026-09-07, turn t0023 rev 4, page
% filed at ../handwritten/20260907-3.10-attempt4-units-of-R-correct.png
% (card 0064, agreed correct on card 0065). Faithful transcription: same
% steps, same order, nothing added.
% Omitted from the transcription, deliberately: rev 4 also carries the answer
% to a hand-check on Z (the pairs a=b=1 and a=b=-1 solving ab=1), which was a
% rung on the way to the answer and not a step of this argument.
% Not supplied by their page, and flagged to them on card 0064:
%   - the intermediate line deg(fg) = deg 1 = 0 with both degrees non-negative;
%     their page goes from additivity straight to deg f = deg g = 0.
% Second question (units in Z_4[x]): the constants, from turn t0023 rev 5
% (card 0065), and the degree-one unit, from t0023 rev 6, page filed at
% ../handwritten/20260907-3.10-attempt6-z4-unit-correct.png (card 0066, agreed
% correct on card 0067). Their earlier claim that 1, 2, 3 are all units there
% was wrong, was sent back on card 0061, and they retracted it themselves.
% The closing description is from turn t0023 rev 7 (card 0067, agreed correct
% on card 0068), transcribed in their own words. One reading was fixed in
% conversation rather than in the text: "all other coefficients are zero
% divisors" must permit 0, or the description fails to cover the constants
% 1 and 3 themselves.
% Not supplied by their pages, and not asked of them:
%   - a proof of the description; card 0067 asked for the description only.
%     (For the record: over Z_4 "zero divisor" and "nilpotent" coincide, and
%     nilpotent is the condition that generalises.)
\begin{solution}
  The units in $R[x]$, for $R$ an integral domain, are
  \[
    \{\, a \in R \subseteq R[x] \;:\; a^{-1} \text{ exists} \,\},
  \]
  the existence of $a^{-1}$ being what is not guaranteed in an integral
  domain.

  They must be constants. For in an integral domain $R$, with $f, g \in R[x]$,
  \[
    \deg(fg) = \deg(f) + \deg(g),
  \]
  since no non-zero coefficients multiply to zero. So $\deg(f) = \deg(g) = 0$.

  In $\mathbb{Z}_{4}[x]$ the units are $1$ and $3$, since $3 \cdot 3 = 9 \equiv
  1$; and also, taking
  \[
    f(x) = 2x + 3, \qquad g(x) = 2x + 3,
  \]
  we have $fg = 4x^{2} + 12x + 9 = 1$.

  So in $\mathbb{Z}_{4}[x]$ the units are those $f$ whose constant term is a
  unit and all of whose other coefficients are zero divisors: constant term in
  $\{1,3\}$, every other coefficient in $\{0,2\}$.
\end{solution}
% ===== END SOLUTION 03.10 =====

\begin{problem}{03.11}
  Show that an element $a$ of an integral domain $R$ is prime if and only if
  $R/(a)$ is an integral domain.
\end{problem}
% ===== SOLUTION 03.11 =====
% PARTIAL. Forward direction only (a prime => R/(a) an integral domain),
% transcribed from the student's slate, 2026-09-07, turn t0026 rev 3, page
% filed at ../handwritten/20260907-3.11-attempt3-forward-correct.png
% (cards 0071, 0072, agreed correct on card 0073). Faithful transcription:
% same steps, same order, nothing added.
% Recorded rather than silently repaired, across the revisions:
%   - rev 1 carried the line "r_1 r_2 is nonzero due to R being an integral
%     domain", which is true but not load-bearing (a class r + (a) is zero
%     when a | r, not when r = 0); flagged on card 0071 and they removed it;
%   - rev 2 justified a nmid 1 with "(a need not be prime to achieve this)",
%     which is false; flagged on card 0072 and they replaced it with the
%     correct hypothesis, "not a unit".
% Not supplied by their page, and not required of them: commutativity of
% R/(a), which is inherited from R.
% Converse direction, divisibility clause only: transcribed from turn t0026
% rev 8, page filed at
% ../handwritten/20260907-3.11-attempt8-converse-divisibility-correct.png
% (cards 0075-0077, agreed correct on card 0078).
% Recorded rather than silently repaired, across those revisions:
%   - revs 5-6 assumed "a nmid bc, a nmid b, a nmid c", which is not the
%     negation of the divisibility clause; the same page then used a | bc.
%     Flagged on cards 0075 and 0076; rev 7 fixed the assumption.
%   - revs 5-7 closed with "and is not a field", which contradicts nothing
%     that was assumed. Flagged on cards 0075 and 0077; rev 8 fixed it.
%   - their page still carries the marker "=>" at the head of this direction,
%     left over from the forward proof; the text below it states the converse,
%     and it is set here as (<=).
% Not supplied by their page, and asked of them on card 0078:
%   - the closing line "contrary to assumption, so a is prime"; the
%     contrapositive stops at "not an integral domain";
%   - the remaining clause of prime, that a is not a unit.
\begin{solution}
  ($\Rightarrow$) Suppose $a \in R$ is prime. Show $R/(a)$ is an integral
  domain.

  Consider
  \[
    (r_{1} + (a))(r_{2} + (a)) = r_{1}r_{2} + (a).
  \]
  Take the factors non-zero, so $a \nmid r_{1}$ and $a \nmid r_{2}$ --- else we
  would have zero divisors. Then $r_{1}r_{2} + (a)$ is non-zero, as
  $a \nmid r_{1}r_{2}$ since $a$ is prime. So $R/(a)$ has no zero divisors.

  Further, $1 + (a) \neq 0 + (a)$, because $a \nmid 1$ ($a$ is not a unit).

  ($\Leftarrow$) Suppose $R/(a)$ is an integral domain. Show $a \in R$ is
  prime.

  Suppose $a$ is not prime. So there exist $b, c \in R$ such that $a \mid bc$,
  but $a \nmid b$ and $a \nmid c$. Note $b + (a)$ and $c + (a)$ lie in
  $R/(a)$, both non-zero since $a \nmid b$, $a \nmid c$. And
  \[
    (b + (a))(c + (a)) = bc + (a) = (a)
  \]
  since $a \mid bc$. So $R/(a)$ has zero divisors and is not an integral
  domain.
\end{solution}
% ===== END SOLUTION 03.11 =====

\begin{problem}{03.12}
  \todo{statement not yet transcribed}
\end{problem}
% ===== SOLUTION 03.12 =====
% TODO(mferguson): your work goes here.
% ===== END SOLUTION 03.12 =====

\begin{problem}{03.14}
  \todo{statement not yet transcribed}
\end{problem}
% ===== SOLUTION 03.14 =====
% TODO(mferguson): your work goes here.
% ===== END SOLUTION 03.14 =====

\begin{problem}{03.15}
  \todo{statement not yet transcribed}
\end{problem}
% ===== SOLUTION 03.15 =====
% TODO(mferguson): your work goes here.
% ===== END SOLUTION 03.15 =====

\begin{problem}{03.17}
  \todo{statement not yet transcribed}
\end{problem}
% ===== SOLUTION 03.17 =====
% TODO(mferguson): your work goes here.
% ===== END SOLUTION 03.17 =====

\begin{problem}{03.25}
  \todo{statement not yet transcribed}
\end{problem}
% ===== SOLUTION 03.25 =====
% TODO(mferguson): your work goes here.
% ===== END SOLUTION 03.25 =====

\begin{problem}{03.26}
  \todo{statement not yet transcribed}
\end{problem}
% ===== SOLUTION 03.26 =====
% TODO(mferguson): your work goes here.
% ===== END SOLUTION 03.26 =====

\end{document}
